% Options for packages loaded elsewhere
\PassOptionsToPackage{unicode}{hyperref}
\PassOptionsToPackage{hyphens}{url}
\documentclass[
]{article}
\usepackage{xcolor}
\usepackage[margin=1in]{geometry}
\usepackage{amsmath,amssymb}
\setcounter{secnumdepth}{-\maxdimen} % remove section numbering
\usepackage{iftex}
\ifPDFTeX
  \usepackage[T1]{fontenc}
  \usepackage[utf8]{inputenc}
  \usepackage{textcomp} % provide euro and other symbols
\else % if luatex or xetex
  \usepackage{unicode-math} % this also loads fontspec
  \defaultfontfeatures{Scale=MatchLowercase}
  \defaultfontfeatures[\rmfamily]{Ligatures=TeX,Scale=1}
\fi
\usepackage{lmodern}
\ifPDFTeX\else
  % xetex/luatex font selection
\fi
% Use upquote if available, for straight quotes in verbatim environments
\IfFileExists{upquote.sty}{\usepackage{upquote}}{}
\IfFileExists{microtype.sty}{% use microtype if available
  \usepackage[]{microtype}
  \UseMicrotypeSet[protrusion]{basicmath} % disable protrusion for tt fonts
}{}
\makeatletter
\@ifundefined{KOMAClassName}{% if non-KOMA class
  \IfFileExists{parskip.sty}{%
    \usepackage{parskip}
  }{% else
    \setlength{\parindent}{0pt}
    \setlength{\parskip}{6pt plus 2pt minus 1pt}}
}{% if KOMA class
  \KOMAoptions{parskip=half}}
\makeatother
\usepackage{color}
\usepackage{fancyvrb}
\newcommand{\VerbBar}{|}
\newcommand{\VERB}{\Verb[commandchars=\\\{\}]}
\DefineVerbatimEnvironment{Highlighting}{Verbatim}{commandchars=\\\{\}}
% Add ',fontsize=\small' for more characters per line
\usepackage{framed}
\definecolor{shadecolor}{RGB}{248,248,248}
\newenvironment{Shaded}{\begin{snugshade}}{\end{snugshade}}
\newcommand{\AlertTok}[1]{\textcolor[rgb]{0.94,0.16,0.16}{#1}}
\newcommand{\AnnotationTok}[1]{\textcolor[rgb]{0.56,0.35,0.01}{\textbf{\textit{#1}}}}
\newcommand{\AttributeTok}[1]{\textcolor[rgb]{0.13,0.29,0.53}{#1}}
\newcommand{\BaseNTok}[1]{\textcolor[rgb]{0.00,0.00,0.81}{#1}}
\newcommand{\BuiltInTok}[1]{#1}
\newcommand{\CharTok}[1]{\textcolor[rgb]{0.31,0.60,0.02}{#1}}
\newcommand{\CommentTok}[1]{\textcolor[rgb]{0.56,0.35,0.01}{\textit{#1}}}
\newcommand{\CommentVarTok}[1]{\textcolor[rgb]{0.56,0.35,0.01}{\textbf{\textit{#1}}}}
\newcommand{\ConstantTok}[1]{\textcolor[rgb]{0.56,0.35,0.01}{#1}}
\newcommand{\ControlFlowTok}[1]{\textcolor[rgb]{0.13,0.29,0.53}{\textbf{#1}}}
\newcommand{\DataTypeTok}[1]{\textcolor[rgb]{0.13,0.29,0.53}{#1}}
\newcommand{\DecValTok}[1]{\textcolor[rgb]{0.00,0.00,0.81}{#1}}
\newcommand{\DocumentationTok}[1]{\textcolor[rgb]{0.56,0.35,0.01}{\textbf{\textit{#1}}}}
\newcommand{\ErrorTok}[1]{\textcolor[rgb]{0.64,0.00,0.00}{\textbf{#1}}}
\newcommand{\ExtensionTok}[1]{#1}
\newcommand{\FloatTok}[1]{\textcolor[rgb]{0.00,0.00,0.81}{#1}}
\newcommand{\FunctionTok}[1]{\textcolor[rgb]{0.13,0.29,0.53}{\textbf{#1}}}
\newcommand{\ImportTok}[1]{#1}
\newcommand{\InformationTok}[1]{\textcolor[rgb]{0.56,0.35,0.01}{\textbf{\textit{#1}}}}
\newcommand{\KeywordTok}[1]{\textcolor[rgb]{0.13,0.29,0.53}{\textbf{#1}}}
\newcommand{\NormalTok}[1]{#1}
\newcommand{\OperatorTok}[1]{\textcolor[rgb]{0.81,0.36,0.00}{\textbf{#1}}}
\newcommand{\OtherTok}[1]{\textcolor[rgb]{0.56,0.35,0.01}{#1}}
\newcommand{\PreprocessorTok}[1]{\textcolor[rgb]{0.56,0.35,0.01}{\textit{#1}}}
\newcommand{\RegionMarkerTok}[1]{#1}
\newcommand{\SpecialCharTok}[1]{\textcolor[rgb]{0.81,0.36,0.00}{\textbf{#1}}}
\newcommand{\SpecialStringTok}[1]{\textcolor[rgb]{0.31,0.60,0.02}{#1}}
\newcommand{\StringTok}[1]{\textcolor[rgb]{0.31,0.60,0.02}{#1}}
\newcommand{\VariableTok}[1]{\textcolor[rgb]{0.00,0.00,0.00}{#1}}
\newcommand{\VerbatimStringTok}[1]{\textcolor[rgb]{0.31,0.60,0.02}{#1}}
\newcommand{\WarningTok}[1]{\textcolor[rgb]{0.56,0.35,0.01}{\textbf{\textit{#1}}}}
\usepackage{graphicx}
\makeatletter
\newsavebox\pandoc@box
\newcommand*\pandocbounded[1]{% scales image to fit in text height/width
  \sbox\pandoc@box{#1}%
  \Gscale@div\@tempa{\textheight}{\dimexpr\ht\pandoc@box+\dp\pandoc@box\relax}%
  \Gscale@div\@tempb{\linewidth}{\wd\pandoc@box}%
  \ifdim\@tempb\p@<\@tempa\p@\let\@tempa\@tempb\fi% select the smaller of both
  \ifdim\@tempa\p@<\p@\scalebox{\@tempa}{\usebox\pandoc@box}%
  \else\usebox{\pandoc@box}%
  \fi%
}
% Set default figure placement to htbp
\def\fps@figure{htbp}
\makeatother
\setlength{\emergencystretch}{3em} % prevent overfull lines
\providecommand{\tightlist}{%
  \setlength{\itemsep}{0pt}\setlength{\parskip}{0pt}}
\usepackage{bookmark}
\IfFileExists{xurl.sty}{\usepackage{xurl}}{} % add URL line breaks if available
\urlstyle{same}
\hypersetup{
  pdftitle={Tidyverse},
  pdfauthor={BIOL 3401 Mount Royal University},
  hidelinks,
  pdfcreator={LaTeX via pandoc}}

\title{Tidyverse}
\author{BIOL 3401 Mount Royal University}
\date{Topic 4, Winter 2026}

\begin{document}
\maketitle

{
\setcounter{tocdepth}{2}
\tableofcontents
}
\section{\texorpdfstring{\textbf{Learning
Objectives}}{Learning Objectives}}\label{learning-objectives}

Tidyverse is a group of packages with functions for organizing tabular
data (data frames/spreadsheet-like data). Many of the functions are
equivalent to operations you performed previously in base R. The whole
point of the Tidyverse is to make your code ``cleaner'' and easier for
humans to understand, especially those of us with minimal programming
experience. It's important for you to know and understand some basic R
functionality, but after today, it's unlikely you will want to go back
to using base R!

After completing this topic, you should feel comfortable:

\begin{enumerate}
\def\labelenumi{\arabic{enumi}.}
\tightlist
\item
  Describing the basic format of of a \texttt{tidyverse} function
\item
  Using some of the key \texttt{tidyverse} functions
\item
  Applying the `pipe' (\%\textgreater\%) operator in \texttt{tidyverse}
\item
  Planning a basic data wrangling task
\item
  Performing a basic data wrangling task using \texttt{tidyverse}
\end{enumerate}

\section{1 - Set up your R workspace for the
day}\label{set-up-your-r-workspace-for-the-day}

\subsection{Step 1a: Set your working
directory}\label{step-1a-set-your-working-directory}

Make sure we're setting our working directory!

\begin{Shaded}
\begin{Highlighting}[]
\CommentTok{\# Set your working directory, then remove "eval = F"!}
\FunctionTok{setwd}\NormalTok{()}
\end{Highlighting}
\end{Shaded}

\subsection{Step 1b: Load packages}\label{step-1b-load-packages}

\begin{Shaded}
\begin{Highlighting}[]
\CommentTok{\# Install the tidyverse package if you haven\textquotesingle{}t already}

\CommentTok{\# Load it}
\FunctionTok{library}\NormalTok{(tidyverse) }
\end{Highlighting}
\end{Shaded}

Now that we have some useful packages loaded, we can start to play with
data.

\subsection{Step 1c: Load your data!}\label{step-1c-load-your-data}

We're going to use a new dataset to explore a number of the
\texttt{tidyverse} functions.

Whenever we work with a new dataset, we need to familiarize ourselves
with what it contains. Let's pretend we don't have metadata for this
dataset (you actually don't), but we \emph{do} have some idea about what
it contains: body measurements from a population of elephants measured
from 1966--1968 and again in 2005--2013.

\begin{Shaded}
\begin{Highlighting}[]
\CommentTok{\# Load the dataset}
\NormalTok{tusks }\OtherTok{\textless{}{-}} \FunctionTok{read\_csv}\NormalTok{(}\StringTok{"datasets/elephant\_tusks.csv"}\NormalTok{)}

\CommentTok{\# Summarize the dataset}
\FunctionTok{summary}\NormalTok{(tusks)}
\end{Highlighting}
\end{Shaded}

Notice that instead of \texttt{read.csv}, we've switched to
\texttt{read\_csv} to load our dataset, which is a tidyverse dataset. It
is useful because it makes what's called a \emph{tibble}, which is a
modern style of tidyverse data frame with some special properties.
You'll see that it's given us a message about the number of rows,
columns, the delimiter (comma-separated, i.e., csv), and its assumptions
about the data type of each column. If you type the name of the object
(tusks) in your R will print out only the first 10 rows of the tibble,
instead of all 780 rows as it would for a basic data frame.

We can check out the class of the tusks tibble:

\begin{Shaded}
\begin{Highlighting}[]
\CommentTok{\# Check the object class}
\FunctionTok{class}\NormalTok{(tusks)}
\end{Highlighting}
\end{Shaded}

R tells us it's a ``spec\_tbl\_df'', ``tbl\_df/tbl'', and
``data.frame''. This means it has a class both as a data frame in base R
and as a tibble for more modern functions.

\section{2 - Planning our data cleaning
task}\label{planning-our-data-cleaning-task}

\subsection{Step 2a: Defining our
endpoint}\label{step-2a-defining-our-endpoint}

By the end of this topic, we need to create a dataset meeting the
following criteria:

\begin{itemize}
\tightlist
\item
  Compares the means and standard deviations of tusk lengths and
  circumferences between study periods
\item
  Includes only adult elephants (\textgreater= 12 years for females and
  \textgreater= 10 years for males)
\item
  Separates female from male elephants
\end{itemize}

\subsection{Step 2b: Understanding the
data}\label{step-2b-understanding-the-data}

The tusks tibble has 780 observations and six columns. Here are the
columns:

\begin{itemize}
\tightlist
\item
  \texttt{sample\_period}, when the observation was made.
\item
  \texttt{sex}, the sex of the elephant (male or female).
\item
  \texttt{estimated\_age}, the age of the elephant in \emph{years}.
\item
  \texttt{shoulder\_height}, the height of the elephant in
  \emph{centimeters}.
\item
  \texttt{tusk\_length}, the length of the tusks in \emph{centimeters}.
\item
  \texttt{tusk\_circumference}, the circumference of the tusks in
  \emph{centimeters}.
\end{itemize}

\subsection{Step 2c: Locate the
obstacles}\label{step-2c-locate-the-obstacles}

\begin{itemize}
\tightlist
\item
  \textbf{Missing data:} A quick glance at the data tells us we have
  many missing values (NA).
\item
  \textbf{Duplicated rows:} You likely won't be able to detect them
  visually because there are so many observations in the data, but there
  are also a few duplicated rows. We will learn how to detect duplicates
  when there are too many rows in the data to locate them visually.
\item
  \textbf{Inconsistencies:} So far, it doesn't look like there are any
  improperly entered values. One clue is that the data types of all the
  rows are as expected. If we had improperly entered measurements, for
  example, the measurement columns (currently numeric) might be
  characters instead.
\end{itemize}

\subsection{Step 2d: Plan decisions}\label{step-2d-plan-decisions}

Starting with the raw data, we are going to execute the following
general steps to produce the endpoint dataset we decided we wanted:

\begin{enumerate}
\def\labelenumi{\arabic{enumi}.}
\tightlist
\item
  Rename some columns for clarity
\item
  Remove the missing and duplicated data
\item
  Add a column that defines the age class of the elephant (juvenile
  versus adult)
\item
  Filter out the juvenile elephants
\item
  Select out the columns we don't need
\item
  Pivot the data into a ``longer'' format to help us summarize
\item
  Summarize by the mean and standard deviation (our final dataset)
\end{enumerate}

We're going to clean the data by running through some functions
belonging to \texttt{tidyr} \& \texttt{dplyr} (two packages that are
automatically loaded when you load \texttt{tidyverse}).

\section{3 - Execute the data cleaning
task}\label{execute-the-data-cleaning-task}

\subsection{\texorpdfstring{Step 3a: \textbf{Change column
names}}{Step 3a: Change column names}}\label{step-3a-change-column-names}

Functions we will use in this step:

\begin{itemize}
\tightlist
\item
  \texttt{rename()}
\end{itemize}

Right now, the measurement units in our data are unclear. Let's first
rename the columns to include units for each of the the four variables
with units (\texttt{estimated\_age}, \texttt{shoulder\_height},
\texttt{tusk\_length}, \texttt{tusk\_circumference}). We will start with
the three columns measured in centimeters.

The \emph{syntax} of using \texttt{rename()} is
\texttt{new\ name\ =\ old\ name}.

\begin{Shaded}
\begin{Highlighting}[]
\CommentTok{\#1 {-} by renaming the columns with a pipe (\%\textgreater{}\%)}
\NormalTok{tusks\_rename }\OtherTok{\textless{}{-}}\NormalTok{ tusks }\SpecialCharTok{\%\textgreater{}\%}
  \FunctionTok{rename}\NormalTok{(}\AttributeTok{estimated\_age\_cm =}\NormalTok{ estimated\_age, }\AttributeTok{shoulder\_height\_cm =}\NormalTok{ shoulder\_height, }\AttributeTok{tusk\_length\_cm =}\NormalTok{ tusk\_length, }\AttributeTok{tusk\_circumference\_cm =}\NormalTok{ tusk\_circumference)}

\CommentTok{\#2 {-} by renaming the columns without a pipe}
\NormalTok{tusks\_rename }\OtherTok{\textless{}{-}} \FunctionTok{rename}\NormalTok{(}\AttributeTok{estimated\_age\_cm =}\NormalTok{ estimated\_age, }\AttributeTok{shoulder\_height\_cm =}\NormalTok{ shoulder\_height, }\AttributeTok{tusk\_length\_cm =}\NormalTok{ tusk\_length, }\AttributeTok{tusk\_circumference\_cm =}\NormalTok{ tusk\_circumference)}
\end{Highlighting}
\end{Shaded}

One of the major new things in how we coded the above block is the pipe
operator (\texttt{\%\textgreater{}\%}). Remember, the pipe passes the
\emph{output} on the left hand side of the pipe as the \emph{input}
expected in the tidyverse function on the next line of code.

This is the basic structure:

\begin{Shaded}
\begin{Highlighting}[]
\CommentTok{\# A tidyverse function without piping}
\NormalTok{output }\OtherTok{\textless{}{-}} \ControlFlowTok{function}\NormalTok{(input, arguments)}
  
\CommentTok{\# The same task using piping}
\NormalTok{output }\OtherTok{\textless{}{-}}\NormalTok{ input }\SpecialCharTok{\%\textgreater{}\%}
  \ControlFlowTok{function}\NormalTok{(arguments)}
\end{Highlighting}
\end{Shaded}

Piping might not seem very useful at the moment wwith only one function,
but we will soon see how useful it can be!

\subsubsection{\texorpdfstring{\textbf{Mini Challenge
1}}{Mini Challenge 1}}\label{mini-challenge-1}

We renamed our measurement columns, but we also need to rename the
estimated\_age column.

\begin{enumerate}
\def\labelenumi{\arabic{enumi}.}
\tightlist
\item
  Rename \texttt{tusks\_rename} to replace \texttt{estimated\_age}
  \texttt{with\ estimated\_age\_yrs}.
\item
  Call the new data \texttt{tusks\_rename\_all}
\end{enumerate}

Use the code block below to code your answers!

\#-------------------------------------------------

\subsection{\texorpdfstring{Step 3b: \textbf{Cleaning the data: Remove
the missing and duplicated
data}}{Step 3b: Cleaning the data: Remove the missing and duplicated data}}\label{step-3b-cleaning-the-data-remove-the-missing-and-duplicated-data}

Functions we will use in this step:

\begin{itemize}
\tightlist
\item
  \texttt{drop\_na()}
\item
  \texttt{arrange()}
\item
  \texttt{distinct()}
\item
  \texttt{group\_by()}
\end{itemize}

We can clearly see the missing data (NA). If we run
\texttt{summary(tusks)}, we can also see the exact number of NAs in each
column at the bottom of the column summary. I also happen to know there
are three duplicate rows in the data. However, you would have a very
difficult time finding them in 780 rows of observations!

\subsubsection{Drop the missing
observations}\label{drop-the-missing-observations}

The first thing we are going to do is drop the missing observations
using the function \texttt{drop\_na()}. We are also going to use the
helper function \texttt{any\_of()} to tell R which columns we want to
remove the missing observations from: \texttt{tusk\_length},
\texttt{tusk\_circumference}, and \texttt{estimated\_age}. This is
because we need these columns for our summary, and missing values will
get in our way. Specifying the columns isn't always necessary, but it's
good to get into the habit of specifying the columns;
\texttt{drop\_na()} removes \emph{all} missing values by default, and
sometimes you might remove rows you need for your analysis just because
there are missing values in another column you don't really care about.
In our case, there is one missing value in the \texttt{shoulder\_height}
column, which we don't really care about, but will remove an observation
with all the other data we \emph{do} care about if we don't specify.

\begin{Shaded}
\begin{Highlighting}[]
\CommentTok{\# Remove NAs, using helper function any\_of() to specify only missing tusk measurements}
\NormalTok{tusks\_no\_na }\OtherTok{\textless{}{-}}\NormalTok{ tusks\_rename }\SpecialCharTok{\%\textgreater{}\%}
  \FunctionTok{drop\_na}\NormalTok{(}\FunctionTok{any\_of}\NormalTok{(}
    \FunctionTok{c}\NormalTok{(}\StringTok{\textquotesingle{}tusk\_length\textquotesingle{}}\NormalTok{, }\StringTok{\textquotesingle{}tusk\_circumference\textquotesingle{}}\NormalTok{))}
\NormalTok{  )}
\end{Highlighting}
\end{Shaded}

\paragraph{\texorpdfstring{\textbf{STOP and check your
work}}{STOP and check your work}}\label{stop-and-check-your-work}

We need to check in frequently to ensure our intermediate data wrangling
steps make sense. We just removed the missing observations from the
data. At this point, we should run the \texttt{summary()} function again
to make sure that we have a dataset with \emph{fewer rows and no NAs in
the tusk\_length\_cm, tusk\_circumference, or estimated\_age columns}.

You should now have a dataset with 597 rows and no NA in either of the
tusk measurement columns.

\subsubsection{Try to find the
duplicates}\label{try-to-find-the-duplicates}

Next, we want to remove all the duplicate rows from the data.

One way to look for duplicates is by using the \texttt{arrange()}
function. This function orders the data according to the values in one
or more columns. This can be an effective way to look for duplicates
because the duplicate rows will now be one after the other.

\subsubsection{\texorpdfstring{\textbf{Mini Challenge
2}}{Mini Challenge 2}}\label{mini-challenge-2}

You can apply the \texttt{arrange()} function using basic tidyverse
syntax or piping.

\begin{enumerate}
\def\labelenumi{\arabic{enumi}.}
\tightlist
\item
  Arrange the \texttt{tusks\_no\_na} dataset by shoulder height
\item
  Call the new data \texttt{tusks\_arrange}
\end{enumerate}

Use the code block below to code your answers!

\#-------------------------------------------------

That didn't work very well. With so few duplicates, and many elephants
having the same shoulder height measurements, it's hard to see the
duplicates.

\subsubsection{Filter out the
duplicates}\label{filter-out-the-duplicates}

Another way we can look for duplicates is by grouping the elephants by
\emph{all} columns and then filtering out the groups with identical data
in all columns. Let's break this down.

\begin{enumerate}
\def\labelenumi{\arabic{enumi}.}
\tightlist
\item
  The \texttt{group\_by()} function groups the data by one or more
  columns. The \emph{syntax} is
  \texttt{group\_by(column\_1,\ column\_2,\ ...\ column\_n)}.
\item
  The \texttt{filter()} function uses relational operators to
  \emph{remove} or \emph{keep} rows that meet our criteria. The
  \emph{syntax} is
  \texttt{column\ name(or\ other\ identifier),\ relational\ operator,\ criteria}.
\item
  The \texttt{n()} function is a helper function to count the number of
  rows in each group (most useful when used with \texttt{group\_by()}).
  This function does not need arguments within the parentheses when used
  within another \texttt{tidyverse} function.
\end{enumerate}

\begin{Shaded}
\begin{Highlighting}[]
\CommentTok{\# Make a new dataset "tusks\_dups" with all duplicate rows to we can see the problem}
\NormalTok{tusks\_dups }\OtherTok{\textless{}{-}}\NormalTok{ tusks\_arrange }\SpecialCharTok{\%\textgreater{}\%}
  \CommentTok{\# Group by all columns in the dataset to pick out 100\% identical rows}
  \FunctionTok{group\_by}\NormalTok{(sample\_period, }
\NormalTok{           sex, }
\NormalTok{           estimated\_age\_yrs, }
\NormalTok{           shoulder\_height\_cm, }
\NormalTok{           tusk\_length\_cm, }
\NormalTok{           tusk\_circumference\_cm) }\SpecialCharTok{\%\textgreater{}\%}
  \CommentTok{\# Filter the groups with more than one row ("groups, greater than, 1")}
  \FunctionTok{filter}\NormalTok{(}\FunctionTok{n}\NormalTok{() }\SpecialCharTok{\textgreater{}} \DecValTok{1}\NormalTok{)}
\end{Highlighting}
\end{Shaded}

The new dataset \texttt{tusks\_dups} has six rows. We can clearly see
that these 6 rows are duplicates of only three distinct rows in the
dataset.

Now that we've assessed the issue, we can remove the duplicates using
the \texttt{distinct()} function. This function's job is to keep only
the rows that are ``distinct'' from other rows in the data, i.e., filter
out the duplicates. It has no arguments other than the dataset name. If
used with a pipe, it can be added on behind the pipe without any
arguments.

\subsubsection{\texorpdfstring{\textbf{Mini Challenge
3}}{Mini Challenge 3}}\label{mini-challenge-3}

\begin{enumerate}
\def\labelenumi{\arabic{enumi}.}
\tightlist
\item
  Use \texttt{distinct()} to remove the duplicate rows from
  \texttt{tusks\_arrange}.
\item
  Call the new dataset \texttt{tusks\_distinct}.
\end{enumerate}

Use the code block below to code your answers!

\paragraph{\texorpdfstring{\textbf{STOP and check your
work}}{STOP and check your work}}\label{stop-and-check-your-work-1}

We just removed three duplicate rows from the data. This means our new
dataset should have 594 rows.

\subsection{\texorpdfstring{Step 3c: \textbf{Enriching the data: Add an
age class
column}}{Step 3c: Enriching the data: Add an age class column}}\label{step-3c-enriching-the-data-add-an-age-class-column}

Functions we will use in this step:

\begin{itemize}
\tightlist
\item
  \texttt{mutate()}
\item
  \texttt{case\_when()}
\end{itemize}

We only want a summary of the adult tusk measurements. Next, we need to
remove the juveniles from the data. We don't have a column for age
class, so we need to make one using the \texttt{mutate()} function.

The \emph{syntax} for \texttt{mutate()} is
\texttt{new\ column\ =\ operator,\ criteria}. This is a simple mutate we
might do if, for example, we wanted to add a column for tusk length in
millimeters (10 mm in 1 cm):

\begin{Shaded}
\begin{Highlighting}[]
\CommentTok{\# Make a new column for tusk length in millimeters}
\FunctionTok{mutate}\NormalTok{(tusks\_distinct, }\AttributeTok{tusk\_length\_mm =}\NormalTok{ tusk\_length\_cm }\SpecialCharTok{*} \DecValTok{10}\NormalTok{)}
\end{Highlighting}
\end{Shaded}

We are going to use relational operators to ask R whether the elephant
in each row is an adult or juvenile, and insert a value in the new age
class column depending on the answer. In fact, we are going to see how
we can give the function multiple relational operators at once. We can
string two relational operators together with two new operators,
\texttt{\textbar{}} and \texttt{\&}. \texttt{\textbar{}} between two
criteria means ``or''. It says, ``find whether the elephant meets
\emph{either} of these criteria. \texttt{\&} between two criteria
means''and''. It says, ``find whether the elephant meets \emph{both} of
these criteria.

In our case, we want ``and'' because we first need to know whether an
elephant is male or female to know whether it would be classified as a
juvenile or adult based on its age (female adult elephants are
\textgreater= 12 years and male adult elephants are \textgreater= 10
years).

The \texttt{case\_when()} function allows us to specify multiple
criteria for a single column. The syntax is
\texttt{if\ criteria\ 1\ matches\ \textasciitilde{}\ insert\ value\ 1,\ if\ criteria\ 2\ matches\ \textasciitilde{}\ insert\ value\ 2,\ ...\ if\ criteria\ n\ matches,\ insert\ value\ n}.
We can specify as many criteria as we want.

Here is an example of how we might separate tall elephants with shoulder
heights \textgreater= 150 cm from short elephants \textless{} 150 cm,
then convert the column to a factor:

\begin{Shaded}
\begin{Highlighting}[]
\CommentTok{\# Add new column for two heights}
\NormalTok{tusks\_distinct }\SpecialCharTok{\%\textgreater{}\%}
  \FunctionTok{mutate}\NormalTok{(}\AttributeTok{split\_height =} \FunctionTok{case\_when}\NormalTok{(shoulder\_height\_cm }\SpecialCharTok{\textgreater{}=} \DecValTok{150} \SpecialCharTok{\textasciitilde{}} \StringTok{"tall"}\NormalTok{,}
\NormalTok{                                  shoulder\_height\_cm }\SpecialCharTok{\textless{}} \DecValTok{150} \SpecialCharTok{\textasciitilde{}} \StringTok{"short"}\NormalTok{)) }\SpecialCharTok{\%\textgreater{}\%}
  \FunctionTok{mutate}\NormalTok{(}\AttributeTok{split\_height\_f =} \FunctionTok{factor}\NormalTok{(split\_height))}
\end{Highlighting}
\end{Shaded}

\subsubsection{\texorpdfstring{\textbf{Mini Challenge
4}}{Mini Challenge 4}}\label{mini-challenge-4}

We need to make a new column with the following criteria to separate
juveniles from adults:

\begin{itemize}
\tightlist
\item
  If sex is female and age is \textgreater= 12, ``ADULT''
\item
  If sex is female and age is \textless{} 12, ``JUVENILE''
\item
  If sex is male and age is \textgreater= 10, ``ADULT''
\item
  If sex is female and age is \textless{} 10, ``JUVENILE''
\end{itemize}

\begin{enumerate}
\def\labelenumi{\arabic{enumi}.}
\tightlist
\item
  Create a new column in the \texttt{tusks\_distinct} dataset called
  \texttt{age\_class}.
\item
  Name the two age classes ``ADULT'' and ``JUVENILE''.
\item
  Create an additional column, \texttt{age\_class\_f}, with the age
  classes as factors.
\item
  Call the new dataset \texttt{tusks\_with\_age\_class}.
\end{enumerate}

Use the code block below to code your answers!

\#-------------------------------------------------

\paragraph{\texorpdfstring{\textbf{STOP and check your
work}}{STOP and check your work}}\label{stop-and-check-your-work-2}

We just added a column with our two age classes as a factor. Our dataset
should now have two categories in the \texttt{age\_class\_f} column,
with 350 adults and 244 juveniles.

\subsection{\texorpdfstring{Step 3d: \textbf{Reshaping the data: Filter
the adults and remove extra
columns}}{Step 3d: Reshaping the data: Filter the adults and remove extra columns}}\label{step-3d-reshaping-the-data-filter-the-adults-and-remove-extra-columns}

Functions we will use in this step:

\begin{itemize}
\tightlist
\item
  \texttt{filter()}
\item
  \texttt{select()}
\end{itemize}

Our next task is to remove the data we don't need so we can focus on
summarizing the data we do need. We need to do the following:

\begin{itemize}
\tightlist
\item
  Remove unnecessary observations (rows)
\item
  Remove unnecessary variables (columns)
\end{itemize}

\subsubsection{Filter out the juveniles (and keep the
adults)}\label{filter-out-the-juveniles-and-keep-the-adults}

We will filter out the juveniles (leaving only the adult observations)
using the \texttt{filter()} function. Filter works something like
indexing and subsetting a data frame by relational operator in base R,
but much more straightforward. The \emph{syntax} of \texttt{filter()} is
\texttt{column\ name,\ relational\ operator,\ criteria}. It looks
something like this:

\begin{Shaded}
\begin{Highlighting}[]
\CommentTok{\# Subset by indexing in base R}
\NormalTok{df[df}\SpecialCharTok{$}\NormalTok{column }\SpecialCharTok{==} \StringTok{"criteria"}\NormalTok{ ,]}

\CommentTok{\# The equivalent, using filter() in tidyverse}
\FunctionTok{filter}\NormalTok{(df, column }\SpecialCharTok{==} \StringTok{"criteria"}\NormalTok{)}

\CommentTok{\# The same thing with piping}
\NormalTok{df }\SpecialCharTok{\%\textgreater{}\%}
  \FunctionTok{filter}\NormalTok{(column }\SpecialCharTok{==} \StringTok{"criteria"}\NormalTok{)}
\end{Highlighting}
\end{Shaded}

\subsubsection{\texorpdfstring{\textbf{Mini Challenge
5}}{Mini Challenge 5}}\label{mini-challenge-5}

\begin{enumerate}
\def\labelenumi{\arabic{enumi}.}
\tightlist
\item
  Filter the \texttt{tusks\_with\_age\_class} data so only adults remain
  (\emph{hint: There is more than one way to do this using filter!})
\item
  Call the new dataset \texttt{tusks\_adults}
\end{enumerate}

Use the code block below to code your answers!

\#-------------------------------------------------

\subsubsection{Select only the columns we
need}\label{select-only-the-columns-we-need}

There are several columns in our data that we no longer need. The only
ones we really need to keep are \texttt{sample\_period}, \texttt{sex},
\texttt{tusk\_length}, and \texttt{tusk\_circumference}. We can safely
get rid of the other variables.

The function that gets rid of columns is \texttt{select()}. It's related
to \texttt{filter()}, but it works on columns instead. We can remove as
many columns as we want up to n-1. The \emph{syntax} for \emph{keeping}
a subset of columns is
\texttt{select(data,\ column1,\ column2,\ ...\ column\_n)}.
\emph{Removing} a subset of columns looks like this:
\texttt{select(data,\ !\ c(column1,\ column2,\ ...\ column\_n))}.

\subsubsection{\texorpdfstring{\textbf{Mini Challenge
6}}{Mini Challenge 6}}\label{mini-challenge-6}

\begin{enumerate}
\def\labelenumi{\arabic{enumi}.}
\tightlist
\item
  Select only the columns from \texttt{tusks\_adults} that we need to
  keep (\texttt{sample\_period}, \texttt{sex},
  \texttt{tusk\_length\_cm}, and \texttt{tusk\_circumference\_cm}).
\item
  Call the new dataset \texttt{tusks\_select}
\end{enumerate}

Use the code block below to code your answers!

\#-------------------------------------------------

\paragraph{\texorpdfstring{\textbf{STOP and check your
work}}{STOP and check your work}}\label{stop-and-check-your-work-3}

We just removed unnecessary columns and filtered out the juvenile
elephants. Our dataset should now have four variables and 350 rows.

\subsection{\texorpdfstring{Step 3e: \textbf{Reshaping the data:
Summarizing tusks by sex and
period}}{Step 3e: Reshaping the data: Summarizing tusks by sex and period}}\label{step-3e-reshaping-the-data-summarizing-tusks-by-sex-and-period}

Functions we will use in this step:

\begin{itemize}
\tightlist
\item
  \texttt{pivot\_longer()}
\item
  \texttt{summarize()}
\item
  \texttt{group\_by()}
\end{itemize}

We are finally ready to summarize the data to meet our goal: the mean
and standard deviation of both tusk measurements, compared between study
periods and separated by sex.

\subsubsection{Pivot the data}\label{pivot-the-data}

Before we summarize, we need to finish one more reshaping task. We
currently have two columns with tusk measurements in our data:
\texttt{tusk\_length\_cm} and \texttt{tusk\_circumference\_cm}. What we
want is the following:

\begin{itemize}
\tightlist
\item
  One column called \texttt{tusk\_dimension} with two factors:
  ``length\_cm'' and ``circumference\_cm''.
\item
  Another column called \texttt{measurement} with the measurement in
  centimeters.
\end{itemize}

We can accomplish this task with the \texttt{pivot\_longer()} function.
\texttt{pivot\_longer()} has several arguments. The necessary ones
include:

\begin{itemize}
\tightlist
\item
  \texttt{cols}, specifying which columns should be ``pivotted'' (moved
  into a single column).
\item
  \texttt{names\_to}, specifying the names of the factor column.
\item
  \texttt{values\_to}, specifying the name of the values column
  corresponding with each factor's observation.
\end{itemize}

We are also going to include another argument in this case:

\begin{itemize}
\tightlist
\item
  \texttt{names\_prefix}, specifying a part of the factor names we want
  to remove.
\end{itemize}

Here is what the pivot looks like:

\begin{Shaded}
\begin{Highlighting}[]
\NormalTok{tusks\_long }\OtherTok{\textless{}{-}}\NormalTok{ tusks\_select }\SpecialCharTok{\%\textgreater{}\%}
  \FunctionTok{pivot\_longer}\NormalTok{(}
    \CommentTok{\# Specifying the columns to pivot}
    \AttributeTok{cols =} \FunctionTok{c}\NormalTok{(tusk\_length, tusk\_circumference),}
    \CommentTok{\# The prefix we want to remove}
    \AttributeTok{names\_prefix =} \StringTok{\textquotesingle{}tusk\_\textquotesingle{}}\NormalTok{,}
    \CommentTok{\# The name of the factor column}
    \AttributeTok{names\_to =} \StringTok{\textquotesingle{}tusk\_dimension\textquotesingle{}}\NormalTok{,}
    \CommentTok{\# The name of the values column}
    \AttributeTok{values\_to =} \StringTok{\textquotesingle{}measurement\textquotesingle{}}
\NormalTok{  )}
\end{Highlighting}
\end{Shaded}

\subsubsection{\texorpdfstring{\textbf{Mini Challenge
7}}{Mini Challenge 7}}\label{mini-challenge-7}

\begin{enumerate}
\def\labelenumi{\arabic{enumi}.}
\tightlist
\item
  Convert the \texttt{tusk\_dimension} column into a factor. \emph{(TIP:
  you can replace an existing column with a different data type using
  the syntax: \texttt{mutate(column\_x\ =\ function(column\_x))} )}
\item
  Call the new dataset \texttt{tusks\_long\_factor}.
\end{enumerate}

Use the code block below to code your answers!

\#-------------------------------------------------

\subsubsection{Summarize the data}\label{summarize-the-data}

To get to our final endpoint with this dataset, we will need to use the
\texttt{summarize()} function. \texttt{summarize()} takes a column and
squishes it down into a single value (mean, sum, standard deviation,
etc.) that summarizes all rows. The \emph{syntax} is:
\texttt{summarize(data,\ summarized\_column\_x\ =\ function(column\_x))}.

There are several useful helper functions we can use within summarize:

\begin{itemize}
\tightlist
\item
  \texttt{mean()}, the mean of all values.
\item
  \texttt{median()}, the median of all values.
\item
  \texttt{sum()}, the sum of all values.
\item
  \texttt{sd()}, the standard deviation of all values.
\end{itemize}

The other important thing to know is that summarize will give us a
single value summarizing all rows \emph{unless} we tell it to summarize
\emph{by group}. We already know we need to do this, because we want a
separate mean for each sex, sample period, and tusk dimension. We
previously used the \texttt{group\_by()} function to group each row by
all variables and then looked for any duplicates. Now we will use
\texttt{group\_by()} to group by all of the variables we want to
summarize by.

\subsubsection{\texorpdfstring{\textbf{Mini Challenge
8}}{Mini Challenge 8}}\label{mini-challenge-8}

Your final challenge is to summarize the mean and standard deviation for
each tusk dimension of each sex in each sample period.

\begin{enumerate}
\def\labelenumi{\arabic{enumi}.}
\tightlist
\item
  Use \texttt{group\_by()} to group \texttt{tusks\_long\_factor} by
  \texttt{sample\_period}, \texttt{sex}, and \texttt{tusk\_dimension}.
  \emph{(NOTE: grouping must happen before summarizing)}
\item
  Summarize the dataset with two new columns called
  \texttt{mean\_measurement} and \texttt{sd\_measurement} for the mean
  and standard deviation of the tusk measurements.
\item
  Call the finalized dataset \texttt{tusks\_cleaned}.
\end{enumerate}

Use the code block below to code your answers!

\#-------------------------------------------------

\paragraph{\texorpdfstring{\textbf{CHECK YOUR FINAL
WORK}}{CHECK YOUR FINAL WORK}}\label{check-your-final-work}

Your final dataset should have:

\begin{itemize}
\tightlist
\item
  8 rows
\item
  5 columns:

  \begin{itemize}
  \tightlist
  \item
    \emph{sample\_period,} a character
  \item
    \emph{sex,}, a character
  \item
    \emph{tusk\_dimension}, a factor
  \item
    \emph{mean\_measurement}, a numeric/double
  \item
    \emph{sd\_measurement}, a numeric/double
  \end{itemize}
\end{itemize}

\end{document}
